% sections/03-methodology.tex — PRISMA multi-stream methodology (spec §5).
\section{Methodology}
\label{sec:method}

\subsection{Study Design}
We conduct a PRISMA-grounded systematization with one deliberate structural change: because the
subject spans academic papers, competition artifacts, vendor disclosures, and registry records,
we maintain \textbf{three evidence streams} rather than forcing all material through a single
review flow. Stream~A (academic literature) follows the classical funnel; Stream~B
(competition/industrial/foundation documents) follows targeted acquisition with quote-level
extraction; Stream~C (registry records, e.g., NVD entries) follows existence verification. The
streams are reported separately and only summed at the top level.

\subsection{Search Strategy}
Stream~A searches used arXiv boolean queries (eight formulations over ``large language model''
/ ``language model'' / ``LLM agent'' crossed with vulnerability discovery, detection,
penetration testing, CTF, exploit, program repair, cyber reasoning system; total 973 raw hits)
and four DBLP queries (214 hits), filtered to January~2024--June~2026 submission dates.
Publisher APIs were unavailable; IEEE/ACM/USENIX coverage came from official indexes and
citation chaining, a limitation we return to below. Streams B/C used targeted retrieval:
official program pages, team/vendor blogs, foundation retrospectives, and NVD. The complete
query log with per-query hit counts is Appendix~B.

\subsection{Selection and Corpus}
Screening operated on a deduplicated pool (316 distinct arXiv IDs plus seed/chaining records).
Inclusion required peer-reviewed status, influential preprint status, or documented-system
evidence relevant to at least one framework cell. The final corpus holds \textbf{36 records}:
23 academic works with full texts converted and validated locally, 2 restricted-access
surveys positioned via verified metadata, 10 competition/industrial primaries with verbatim
quote extraction, and 1 registry record. Four candidate records were excluded (wrong
identifier, press-only claims, out-of-window scope, unresolvable primary URL). Saturation was
tested non-circularly across six novelty axes (new systems, new loop mechanisms, new validity
problems, contradictions, new industrial evidence, negative results); two consecutive passes
without axis-level novelty stopped inclusion.

\subsection{Quality Rubric and Evidence Classes}
Each record carries reproducibility, baseline-rigor, and artifact-availability scores (0--3),
plus evidence class (benchmark, controlled experiment, production-network experiment,
competition, production discovery, incident report, measurement, registry), unit type, access
status, and provenance type (PDF+TXT, HTML-extract, metadata-only). Vendor-reported statistics
are labeled CLAIMED regardless of source authority; third-party reproduction passes are marked
as such.

\subsection{Classification Protocol}
Autonomy (A0--A3) is assigned by the decision-authority test --- who chooses the next action? ---
with A3 additionally requiring independent planning, execution, validation, and reporting in
environment class E2 (production or production-like external targets; E0=synthetic,
E1=realistic-sandboxed). Claims are graded by their weakest missing condition. Loop position
uses the set $\{$B1..B4$\}$, recorded as sets for multi-phase systems. Every assignment's
behavioral evidence is published in the artifact (\texttt{autonomy-loop-assignments.csv});
benchmarks carry probed phases rather than axis positions, keeping thematic coverage out of the
taxonomy.

\subsection{Limitations}
First-page screening caps per query, publisher-API absence, single-team classification without
inter-rater statistics, and English-language sources are documented limitations; none, we argue,
plausibly reverses the central negative finding (Sec.~\ref{sec:conclusion}), which is robust to
reasonable variations of the inclusion rule.
